\section{Scalars and Vectors}
\label{sec:scalars-vectors}

\subsection{Why Direction Matters}

A numerical value does not always provide a complete physical description. The
statement
\[
v=20\,\mathrm{m\,s^{-1}}
\]
specifies a speed, but it does not say whether the object moves east, west,
upward, or in some other direction. By contrast,
\[
\mathbf v=20\,\mathrm{m\,s^{-1}}\,\mathbf e_x
\]
specifies both the magnitude and the direction of the motion. This distinction
separates \emph{scalar} quantities from \emph{vector} quantities.

\begin{bookinsight}[Two kinds of physical information]
A scalar answers ``how much?'' A vector answers both ``how much?'' and ``in
which direction?''
\end{bookinsight}

\subsection{Scalars}

A scalar is completely characterized by a numerical value and a unit. Typical
examples are mass, temperature, electric potential, pressure, energy, and
probability density. For example,
\[
T=300\,\mathrm{K},
\qquad
m_e=9.11\times10^{-31}\,\mathrm{kg}.
\]
A scalar may vary from point to point or with time, but it does not acquire a
spatial direction. Temperature in a room, for instance, can be represented by
a scalar field $T(\mathbf r,t)$.

\subsection{Vectors}

A vector possesses a magnitude and a direction. Displacement, velocity,
acceleration, force, momentum, electric field, magnetic field, and probability
current are vector quantities. We write vectors in boldface,
\[
\mathbf A,
\]
and reserve ordinary italic symbols such as $A$ for scalars or for the
magnitude $A=\lVert\mathbf A\rVert$ when the meaning is clear.

In a Cartesian basis, a vector is represented by components:
\[
\mathbf A=A_x\mathbf e_x+A_y\mathbf e_y+A_z\mathbf e_z.
\]
The components depend on the selected coordinate system, whereas the geometric
vector does not. Rotating the axes changes the numbers $A_x,A_y,A_z$ but does
not change the physical displacement, velocity, or force represented by
$\mathbf A$.

\begin{bookfigure}{A vector resolved into Cartesian components. The diagonal
arrow is the geometric vector; the horizontal and vertical arrows are its
components in the chosen basis.}
\begin{tikzpicture}[scale=1.05]
  \draw[physics axis] (-0.2,0) -- (5.6,0) node[right] {$x$};
  \draw[physics axis] (0,-0.2) -- (0,4.2) node[above] {$y$};
  \draw[physics vector] (0,0) -- (4.4,3.15)
    node[above right,visual label] {$\mathbf A$};
  \draw[construction line] (4.4,0) -- (4.4,3.15);
  \draw[construction line] (0,3.15) -- (4.4,3.15);
  \draw[physics vector secondary] (0,-0.25) -- (4.4,-0.25)
    node[midway,below] {$A_x$};
  \draw[physics vector secondary] (-0.25,0) -- (-0.25,3.15)
    node[midway,left] {$A_y$};
  \draw[visualgold,line width=0.9pt,-{Latex[length=1.5mm]}]
    (0.9,0) arc[start angle=0,end angle=35.6,radius=0.9];
  \node[text=visualgold,font=\small] at (1.12,0.36) {$\theta$};
\end{tikzpicture}

\end{bookfigure}

\begin{visualguidebox}
The components are not separate vectors hidden inside $\mathbf A$. They are the
projections of the same vector onto the selected basis directions.
\end{visualguidebox}

\subsection{Transformation Test}

A useful way to distinguish scalars from vectors is to ask how the quantity
changes when the coordinate axes are rotated. A true scalar is unchanged. A
vector keeps its magnitude and geometric direction, but its components
transform together according to the rotation. This transformation law, rather
than the presence of an arrow in a diagram, is the precise mathematical
criterion for vector character.

\begin{warningbox}
A quantity is not a vector merely because it has several numerical entries.
Those entries must transform as vector components under a change of basis.
\end{warningbox}
